% 颜色 key 别名回归测试（issue #163）
%
% 每个模块的「文字颜色」历史上有两种拼法：core 层用 font-color / 字体颜色，
% 批注/眉批/侧批/印章用 color / 颜色。写错的那一种以前被 unknown-key
% fallback 静默吞掉（#163：\夹注设置{color=red} 完全无效）。
%
% 本文件让每个模块用一个独一无二的颜色，且刻意混用两种拼法：
% 只要有一个别名失效，对应颜色就不会出现在 PDF 内容流里。
% test/color_test.py 直接对内容流做断言，不依赖基线图像。
\documentclass{ltc-guji}
\setmainfont{TW-Kai}
\关闭分页
\开启句读模式

% 夹注：全局用中文 `颜色`，局部用英文 `color`（两者都是 font-color 的别名）
\夹注设置{颜色={255, 0, 0}}
% 侧批：模块原生是 `color`，这里刻意用 core 层拼法 font-color
\侧批设置{font-color={0, 128, 0}}
% 批注：刻意用 `字体颜色`（原生是 `颜色`）
\批注设置{字体颜色={255, 0, 255}}
% 眉批：同上，用 font-color
\眉批设置{font-color={0, 128, 128}}
% 句读：原生是 judou-color / 句读颜色，这里用裸 `颜色`
\句读设置{颜色={255, 128, 0}}

\begin{document}
\begin{正文}
一二三，四五六。

% 夹注全局色（红）+ 局部覆盖（蓝，英文 color）
天地\夹注{全局红色夹注}玄黄\夹注[color={0, 0, 255}]{局部蓝色}宇宙。

% 侧批（绿）
日月\侧批{侧批绿色}盈昃。

% 文本框：core 层原生 font-color，这里用裸 color
\文本框[color={128, 0, 128}, font-size=20pt]{文本框紫色}

% 行（Column）：原生 font-color，这里用中文 `颜色`
\行[颜色={0, 0, 128}]{整列藏青}

\批注[x=100pt, y=120pt, height=4]{批注品红}
\眉批[height=3, x=260pt]{眉批靛青}

\end{正文}
\end{document}
